\documentclass[a4paper]{scrreprt} % Рекомендований клас для навчальних видань

% lecture_onu.sty
% Стильовий файл, адаптований до вимог ОНУ (14 кегель, 1.5 інтервал, поля 2 см)
\ProvidesPackage{lecture_onu}

% !TeX program = xelatex
% !TeX options = -shell-escape
\usepackage{ragged2e}
\usepackage{cmap}
\usepackage[ukrainian]{babel}

% 1. МОВА: Встановлено українську мову
\usepackage{algorithm}
\usepackage{algorithmic}


%tabularx, бо він сам розподіляє простір і не дає таблиці «вилазити» за поля.
\usepackage{tabularx}

% 3, 4. ПОЛЯ ТА РОЗМІР: Встановлено поля 2 см та розмір шрифту 14 пт (для KOMA-script)
\KOMAoptions{fontsize=14pt}
\usepackage[left=2cm, right=2cm, top=2cm, bottom=2cm]{geometry}


\usepackage{iftex}
\ifPDFTeX
\usepackage[
activate={true,nocompatibility},
final,
tracking=true,
factor=1100,
stretch=10,
shrink=10,
letterspace=700
]{microtype}
\else
\usepackage[stretch=10,shrink=10,tracking=true]{microtype}
\fi


% 2. ШРИФТ: Видалено kpfonts. Встановлюємо Times New Roman.
\usepackage{fontspec}
% ПРИМІТКА: Переконайтеся, що "Times New Roman" встановлений у вашій системі. 
% Для Linux-систем може знадобитися Tinos або Liberation Serif.
\setmainfont{Times New Roman}
\setsansfont{Times New Roman}
\setmonofont{Courier New} % Змінено на більш стандартний моноширинний

%автоматично коригує інтервали між буквами та словами
\usepackage{microtype}


% Видалено mathastext та пов'язані пакети kpfonts для дотримання формату
\usepackage{amsmath,amssymb,amsfonts}
\usepackage{icomma}
\usepackage{mathtools}
\usepackage{nicefrac}
\usepackage{amsthm}


% 5. ІНТЕРВАЛ ТА ВІДСТУП: Встановлено 1.5 інтервал. Видалено parskip.
\usepackage{setspace}
\onehalfspacing 
\setlength{\parskip}{0pt} % Видалено додатковий інтервал між абзацами
\setlength{\parindent}{1cm} % Залишено стандартний абзацний відступ

\usepackage{enumitem}
\setlist{noitemsep} % Залишено, для охайних списків
\renewcommand{\labelitemi}{\faCheckSquareO}

\usepackage[bottom]{footmisc}
\setlength{\footnotesep}{12pt}%

\usepackage[
automark,
headsepline,
footsepline,
autooneside
]{scrlayer-scrpage}
\renewcommand*{\chaptermarkformat}{}
\renewcommand*{\sectionmarkformat}{}
\automark[]{chapter}
\ihead[]{\leftmark}
\chead[]{}
\ohead[]{}




% 7. КОПІРАЙТ: Змінено на українську мову
\ifoot[\textcopyright~{Платонова Є. В, Страхов Є. М.}]{\textcopyright~{ Платонова Є. В, Страхов Є. М.}}
\cfoot[]{}
\ofoot[\pagemark]{\pagemark}
\setkomafont{pageheadfoot}{\footnotesize}
\setkomafont{pagehead}{\small\bfseries\itshape}
\setkomafont{pagenumber}{\normalsize}
\pagestyle{scrheadings}

% СТИЛІЗАЦІЯ ЗАГОЛОВКІВ (Залишено, але пам'ятайте, що вона не є стандартною для ДСТУ)
\usepackage{titlesec}
\titleformat{\chapter}[display]
{\sffamily\bfseries}
{\begin{tikzpicture}[remember picture,overlay]
		\node[below left, shade, top color=black, bottom color=pyblue, rectangle, inner sep=2pt, xshift=-2.0cm, minimum width=1.75cm, minimum height=5cm, text width=1.75cm, align=center] at (current page.north east) {\color{pyyellow}\fontsize{50}{60}\selectfont\mbox{}\\[2cm]\thechapter};
\end{tikzpicture}}
{-1.425cm}
{\Huge}

\titlespacing*{\chapter}
{0pt}{0pt}{30pt}

\usepackage{fancyvrb}
\usepackage[dvipsnames,svgnames,x11names]{xcolor}
\usepackage{minted}
\usemintedstyle{vs}
\usepackage[many]{tcolorbox}
\tcbuselibrary{skins}
\tcbuselibrary{minted}
\tcbuselibrary{theorems}

\titleformat{\subsection}[block] % block = окремий рядок
{\bfseries\normalsize}         % жирний + стандартний розмір
{}                             % номер (порожньо для subsection*)
{0pt}                          % відступ
{}   



% Оформлення блоків (прикладів, задач) залишено.
\newtcolorbox[auto counter, number within=chapter]{pr}[1]{
	enhanced, 
	arc=0mm, outer arc=0mm,
	frame hidden, 
	borderline north={1pt}{-1pt}{MidnightBlue}, borderline south={2pt}{-2pt}{MidnightBlue},
	boxrule=0pt, titlerule=0pt,
	left=0.2cm, right=0.2cm,
	toptitle=0.1cm, bottomtitle=0.1cm, top=0.2cm, colbacktitle=MidnightBlue,
	colback=GhostWhite, coltitle=white, before skip=12pt plus 2pt minus 2pt, after skip=10pt plus 2pt minus 2pt, fonttitle=\bfseries, title={Задача~\thetcbcounter}, #1
} % умови задач

\newtcblisting{py}[1]{
	enhanced, arc=0pt, outer arc=0pt, 
	listing engine=minted, minted style=vs, minted language=python, minted options={fontsize=\small, numbersep=3mm, baselinestretch=1.1}, fonttitle=\bfseries, title={#1}, 
	listing only,
	colback=MintCream, colbacktitle=pyblue,
	frame hidden, 
	borderline north={1pt}{-1pt}{pyblue}, borderline south={2pt}{-2pt}{pyblue}, 
	boxrule=0pt, titlerule=0pt,
	left=0.2cm, right=0.2cm,  
	toptitle=0.1cm, bottomtitle=0.1cm, top=0.2cm, before skip=10pt plus 2pt minus 2pt, after skip=10pt plus 2pt minus 2pt, after={\par\noindent}
} % код

\newtcblisting{pydef}{
	enhanced, arc=0pt, outer arc=0pt, 
	listing engine=minted, minted style=vs, minted language=python, minted options={fontsize=\small, numbersep=3mm}, 
	colback=orange!10!white,
	listing only,
	frame hidden, 
	borderline north={1pt}{-1pt}{orange}, borderline south={2pt}{-2pt}{orange}, 
	boxrule=0pt, titlerule=0pt,
	left=0.2cm, right=0.2cm, 
	toptitle=0.1cm, bottomtitle=0.1cm, top=0.2cm, 
	before skip=10pt plus 2pt minus 2pt, after skip=10pt plus 2pt minus 2pt, after={\par\noindent}
} % шпаргалки с кодом

\newtcolorbox{defn}[1]{
	enhanced, arc=0pt, outer arc=0pt, 
	colback=orange!10!white,
	fonttitle=\bfseries, title={#1}, before skip=10pt plus 2pt minus 2pt, after skip=10pt plus 2pt minus 2pt,
	colbacktitle=orange,  
	frame hidden, 
	borderline north={1pt}{-1pt}{orange}, borderline south={2pt}{-2pt}{orange}, 
	boxrule=0pt, titlerule=0pt,
	left=0.2cm, right=0.2cm,  toptitle=0.1cm, bottomtitle=0.1cm, top=0.2cm
} % головне

\newtcolorbox{pyr}{
	enhanced, arc=0pt, outer arc=0pt, 
	colback=white, 
	frame hidden, 
	before skip=10pt plus 2pt minus 2pt, after skip=10pt plus 2pt minus 2pt, 
	boxrule=0pt, 
	left=-0.1cm, right=-0.1cm, titlerule=0em, toptitle=0cm, bottomtitle=0cm, top=0.2cm
} % код + результаты

\newtcolorbox{fl}[1]{
	arc=0pt, outer arc=0pt, colback=Gainsboro!30, colframe=darkgray, left=2mm, enhanced, coltitle=white, colbacktitle=black, 
	fonttitle=\ttfamily\bfseries, title={#1}, 
	frame hidden, 
	borderline north={1pt}{-1pt}{black}, borderline south={2pt}{-2pt}{black},
	boxrule=0pt, titlerule=0pt,
	left=0.2cm, right=0.2cm,  toptitle=0.1cm, bottomtitle=0.1cm, top=0.2cm, fontupper=\small\ttfamily
} % файли

\newtcolorbox{res}[1]{
	enhanced, arc=0pt, outer arc=0pt, 
	fonttitle=\bfseries, title={#1}, fontupper=\small, colback=MintCream, colframe=MediumBlue, listing only, 
	boxrule=0pt, titlerule=0pt,
	frame hidden, 
	borderline north={1pt}{-1pt}{MediumBlue}, borderline south={2pt}{-2pt}{MediumBlue},
	left=0.2cm, right=0.2cm,  toptitle=0.1cm, bottomtitle=0.1cm, top=0.2cm, 
	before skip=10pt plus 2pt minus 2pt, after skip=10pt plus 2pt minus 2pt
} % результати

% ... решта пакетів та макросів залишаються, якщо вони не суперечать вимогам
\usepackage{graphicx}
\usepackage{tikz}
\usepackage{pgfplots}




\usepackage{hyperref}
\hypersetup{
	pdftitle={Конспект лекцій}, % Змінено заголовок
	pdfauthor={Євген Страхов}, % Змінено автора
	unicode=true,
	colorlinks=true,
	bookmarks=true,
	linkcolor=MediumBlue,
	citecolor=FireBrick,
	urlcolor=DodgerBlue
}

\usepackage{menukeys}
\usepackage[titles]{tocloft}
\renewcommand{\cftchappagefont}{\rmfamily}
\renewcommand{\cftpartpagefont}{\rmfamily}
\setcounter{secnumdepth}{2}

\newcommand{\mi}[1]{\mintinline{python}{#1}}
\usepackage{python}
\usetikzlibrary{fadings}
\newcommand\fadingtext[3][]{%
	\begin{tikzfadingfrompicture}[name=fading letter]
		\node[text=transparent!0,inner xsep=0pt,outer xsep=0pt,#1] {#3};
	\end{tikzfadingfrompicture}%
	\begin{tikzpicture}[baseline=(textnode.base)]
		\node[inner sep=0pt,outer sep=0pt,#1](textnode){\phantom{#3}}; 
		\shade[path fading=fading letter,#2,fit fading=false]
		(textnode.south west) rectangle (textnode.north east);% 
	\end{tikzpicture}% 
}
\definecolor{pyblue}{HTML}{3C7DA7}
\definecolor{pyyellow}{HTML}{FFDB6D}
\newcommand{\Python}{{\fontspec[Scale=1.17]{Inconsolata}\color{pyblue}python}}
\usepackage{fontawesome}
\newcommand{\myhref}[2]{\href{#1}{\sffamily #2~\faExternalLink}}
\usepackage{booktabs}
\usepackage{makecell}
\usepackage{multirow}
\usepackage{multicol}
\usepackage{longtable}
\RequirePackage[square,numbers,sort&compress]{natbib}
\setlength{\bibsep}{6pt}

% 9. ЗМІСТ: Встановлено українську назву
\addto\captionsukrainian{%
	\renewcommand{\contentsname}{Зміст}%
}

% 8. ТИТУЛЬНИЙ АРКУШ: Змінено текст на український
\newcommand{\mytitlepage}[2]{
	% \ThisTileWallPaper... (Зображення фону, можливо, доведеться видалити для формальності)
	\thispagestyle{empty}
	\vspace*{2cm}
	\begin{tcolorbox}[
		beamer,
		interior style={left color=LemonChiffon,right color=yellow},
		boxrule=0pt,
		colframe=white,
		arc=10pt,
		outer arc=10pt,
		remember as=title,
		leftrule=0.25cm,
		rightrule=0pt,
		fontupper=\color{white}\sffamily\bfseries\huge,
		enlarge left by=-1.2in-\hoffset-\oddsidemargin, 
		width=0.7\paperwidth, 
		left=1cm,
		top=0.3cm, 
		bottom=0.3cm, 
		]
		{\large\bfseries\color{MidnightBlue}\fontspec[Ligatures=TeX]{Exo 2 Semi Bold Expanded}
			Основи програмування. Курс лекцій}
	\end{tcolorbox}
	\vspace{8cm}
	\begin{tcolorbox}[
		beamer,
		interior style={left color=MidnightBlue,right color=DodgerBlue},
		boxrule=3pt,
		colframe=white,
		arc=10pt,
		remember as=title,
		leftrule=3pt,
		rightrule=0pt,
		fontupper=\color{white}\sffamily\bfseries\huge,
		enlarge left by=-\hoffset-\oddsidemargin, 
		enlarge right by=-\paperwidth+1in+\hoffset+\oddsidemargin+\textwidth,
		width=0.9\paperwidth, 
		left=1cm, 
		right=\paperwidth-1in-\hoffset-\oddsidemargin-\textwidth,
		top=0.75cm, 
		bottom=0.75cm, 
		]
		{\noindent\hfill\huge\bfseries\fontspec{Exo 2 Semi Bold Expanded}
			Частина~#1\\[0.1cm]
			\mbox{}\hfill\MakeUppercase{#2}}
	\end{tcolorbox}
	\vfill
	\begin{tcolorbox}[
		beamer,
		interior style={left color=LemonChiffon,right color=yellow},
		boxrule=0pt,
		colframe=white,
		arc=10pt,
		outer arc=10pt,
		remember as=title,
		leftrule=0.25cm,
		rightrule=0pt,
		fontupper=\color{white}\sffamily\bfseries\huge,
		enlarge left by=-1.2in-\hoffset-\oddsidemargin, 
		width=0.7\paperwidth, 
		left=1cm,
		top=0.3cm, 
		bottom=0.3cm, 
		]
		{\normalsize\bfseries\color{MidnightBlue}\fontspec{Exo 2 Semi Bold Expanded}
			Євген Страхов\\[-5\jot]
			Одеський національний університет, 2017} % Змінено рік
	\end{tcolorbox}
}


% ---------- Packages ----------
\usepackage{wrapfig}     % figures with text wrapped around them
\usepackage{textcomp}    % extra text symbols
\usepackage{mathrsfs}    % \mathscr
\usepackage{array}       % needed for \newcolumntype

% ---------- TikZ libraries ----------
\usetikzlibrary{arrows}
\usetikzlibrary{arrows.meta}
\usetikzlibrary{shadows}
\usetikzlibrary{shadows.blur}
\usetikzlibrary{shapes}


% ---------- Tables ----------
\renewcommand{\arraystretch}{1.1}

% Format a whole table row: use $ for the first column and ^ for the others, then \rowstyle{\bfseries}
\newcolumntype{$}{>{\global\let\currentrowstyle\relax}}
\newcolumntype{^}{>{\currentrowstyle}}
\newcommand{\rowstyle}[1]{\gdef\currentrowstyle{#1}%
	#1\ignorespaces
}

% ---------- Commands ----------
% Lab header with the Python logo (needs Py.png)
\newcommand{\myheader}[1]{
	\begin{tcolorbox}[enhanced, fonttitle=\scriptsize\sffamily\bfseries, coltitle=pyyellow, colframe=white, title style={left color=pyblue, right color=black}, colback=white, boxrule=0pt, leftrule=0pt, rightrule=0pt, titlerule=0pt, toprule=0pt, bottomrule=1pt, arc=0pt, outer arc=0pt, left=0mm, right=0mm, adjusted title={Курс <<Основи програмування>>\hbox to 1.75cm {}\\Спеціальність <<Математика>>\hbox to 1.75cm {} }, halign title=flush right, boxsep=0mm, toptitle=1.5mm, bottomtitle=1.5mm]
		\begin{minipage}[b]{0.75\textwidth}
			{\Large\sffamily\bfseries #1}
			\vspace{-0.3cm}
		\end{minipage}
		\hfill
		\begin{minipage}{0.15\textwidth}
			\vspace{-0.65cm}
			\includegraphics[width=\linewidth]{Py.png}
			\vspace{-0.6cm}
		\end{minipage}
		\par\vspace{6pt}
		\begin{tikzpicture}
			\fill[left color=pyblue, right color=black] (0,0) rectangle (\textwidth,-1pt);
		\end{tikzpicture}
	\end{tcolorbox}
	\bigskip
}


% Визначення кольорів градієнта
\definecolor{labdarkblue}{RGB}{20, 50, 80}
\definecolor{lablightblue}{RGB}{45, 110, 160}

% Створення команди \LabHeader{Лабораторна}{Варіант}
\newcommand{\LabHeader}[2]{%
	\noindent
	\begin{tikzpicture}[remember picture]
		% 1. Градієнтна смуга зверху
		\shade[left color=lablightblue, right color=black!90] 
		(0,0) rectangle (\textwidth, 0.7cm);
		
		% 2. Текст "Лабораторна X" білим кольором справа на смузі
		\node[anchor=east, pyyellow, font=\bfseries\sffamily\large] 
		at (\textwidth - 1.5cm, 0.35cm) {Лабораторна #1};
		
		% 3. Текст "Варіант Y" знизу ліворуч
		\node[anchor=west, black, font=\bfseries\Large] 
		at (0, -0.6cm) {Варіант #2};
		
		% 4. Зображення Python (з виходом за межі смуги)
		% Переконайтеся, що файл Py.png існує в папці проєкту
		\node[anchor=north east] 
		at (\textwidth + 0.2cm, 0.8cm) {\includegraphics[width=1.8cm]{Py.png}};
		
		% 5. Нижня роздільна лінія
		\draw[line width=0.8pt, lablightblue] 
		(0, -1.1cm) -- (\textwidth, -1.1cm);
	\end{tikzpicture}%
	\vspace{0.5cm} % Відступ після шапки
}



% Variant heading (Варіант = Ukrainian; pylab.sty had Russian "Вариант")
\newcommand{\variant}[1]{
	\noindent\textsf{\large\bfseries Варіант~#1}
}

% Link that opens in a new window, with an icon
\newcommand{\myurl}[2]{
	\href[pdfnewwindow=true]{#1}{#2~\faExternalLink}
}

% Inline Python code.
% NOTE: pylab.sty calls this \py, but main.tex already has an environment
% named "py" (\newtcblisting{py}), which defines \py. So we call it \pyi.
% main.tex's \mi does the same thing.
\newcommand{\pyi}[1]{\mintinline{python}{#1}}

% Number or letter in a filled circle
\newcommand*\circled[1]{\tikz[baseline=(char.base)]{
		\node[shape=circle,draw=pyblue, fill=pyblue,inner sep=2pt] (char) {\color{pyyellow}\bfseries\sffamily#1};}}

% ---------- Lists ----------
% enumerate items get circled letters: (A), (B), ...
\setlist[enumerate]{
	label=\protect\circled{\Alph*},
	itemsep=6pt
}



\begin{document}
	
	% === ТИТУЛЬНА СТОРІНКА ===
	% Необхідно заповнити фактичними даними:
	% \mytitlepage{Частина №}{Підзаголовок}
	%\mytitlepage{1}{Вступ до Програмування} 
	
\titleformat{\section}{\normalfont\Large\bfseries}{\thesection}{1em}{}
	
	
	
	
	% === ОСНОВНИЙ ТЕКСТ ЛЕКЦІЇ ===
	% Вміст із файлу l1.tex:
		%----------------------------------------------------------------------------------------
%	TITLE PAGE
%----------------------------------------------------------------------------------------
\KOMAoptions{fontsize=12pt} 
\begin{titlepage} % Suppresses displaying the page number on the title page and the subsequent page counts as page 1
	\newcommand{\HRule}{\rule{\linewidth}{0.5mm}} % Defines a new command for horizontal lines, change thickness here
	
	\centering % Centre everything on the page
	
	\Large МІНІСТЕРСТВО ОСВІТИ І НАУКИ УКРАЇНИ\\  % Main heading such as the name of your university/college
	\medskip
	\large ОДЕСЬКИЙ НАЦІОНАЛЬНИЙ УНІВЕРСИТЕТ ІМЕНІ І.І. МЕЧНИКОВА% Main heading such as the name of your university/college
	
	\large ФАКУЛЬТЕТ МАТЕМАТИКИ, ФІЗИКИ ТА ІНФОРМАЦІЙНИХ ТЕХНОЛОГІЙ % Major heading such as course name
	
	%\textbf{\Large Кафедра Оптимального Керування та Економічної Кібернетики}\\[1.5cm] % Minor heading such as course title
	
	\vspace*{5\baselineskip} % Whitespace above the title
	
	
	{\Large 	\textbf{Євгенія ПЛАТОНОВА, Євген СТРАХОВ}\\} % Editor list
	
	\vspace{1\baselineskip} % Whitespace below the editor list
	\rule{\textwidth}{1.6pt}\vspace*{-\baselineskip}\vspace*{2pt} % Thick horizontal rule
	\rule{\textwidth}{0.4pt} % Thin horizontal rule
	
	\vspace{1\baselineskip} % Whitespace above the title
	
	{\LARGE \textbf{Програмування}} % Title
	
	\vspace{0.5\baselineskip} % Whitespace below the title
	
	\rule{\textwidth}{0.4pt}\vspace*{-\baselineskip}\vspace{3.2pt} % Thin horizontal rule
	\rule{\textwidth}{1.6pt} % Thick horizontal rule
	
	\vspace{1\baselineskip} % Whitespace after the title block
	
	%------------------------------------------------
	%	Subtitle
	%------------------------------------------------
	
	\LARGE \textbf{Основи програмування на мові Python.\\ Структурне програмування\\} % Subtitle or further description
	%{(адаптація learnopengl.com)\\}
	
	\vspace*{0.5\baselineskip} % Whitespace under the subtitle
	
	\Large \textbf{Практикум. Збірник завдань до лабораторних робіт }
	\vspace*{1.5\baselineskip}
	%\vfill\v% Position the date 3/4 down the remaining page
	
	%		Рекомендовано Методичною радою ОНУ імені І.І. Мечникова\\
	%		як навчальний посібник для здобувачів ступеня бакалавра\\ 
	%		спеціальності 113 «Прикладна математика», 122 «Комп’ютерні науки»,\\
	%		123 «Комп’ютерна інженерія», 126 «Інформаційні системи та технології»\\
	
	
	\vspace*{0.25\baselineskip}
	
	%	\texttt{Електронне мережне навчальне видання}
	\texttt{Електронне навчальне видання}
	
	%{\large\today} % Date, change the \today to a set date if you want to be precise
	
	%------------------------------------------------
	%	Logo
	%------------------------------------------------
	
	%------------------------------------------------
	%	Publisher
	%------------------------------------------------
	\vspace*{4\baselineskip}
	{ \large 
		ОДЕСА\\ \vspace{-1.5mm}
		ОНУ\\ \vspace{-2.5mm}
		2026 }%Publication year
	%\vfill\vfill
	%\includegraphics[width=0.2\textwidth]{placeholder.jpg}\\[1cm] % Include a department/university logo - this will require the graphicx package
	
	%----------------------------------------------------------------------------------------
	
	
	\pagebreak
	
	\large
	\raggedright \Large {\textbf{УДК 004.4:004.421 }\\
		\hspace*{1 cm} XXX}
	
	\bigskip
	
	\large
	\hspace{1cm}\textbf{Автори:}\\
	\textbf{\textit{Є.В. Платонова}} старший викладач кафедри оптимального керування і економічної кібернетики.\\ 
	\textbf{\textit{Є.М. Страхов}} канд. фіз.-мат. наук, доцент, доцент кафедри оптимального керування і економічної кібернетики.\\ 
	\medskip
	\hspace{1cm}\textbf{Рецензенти:}\\
	\textbf{\textit{В.О. Сперанський }}  канд. техн. наук, доцент, доцент кафедри комп’ютеризовані системи та програмні технології НУ "Одеська Політехніка";\\
	\textbf{\textit{}} ;\\
	\textbf{\textit{,}} .
	%		\begin{tabbing}
		%			\textbf{Рецензенти: }  \= \hspace{ 1cm} \= \kill
		%			\textbf{Рецензенти:} \> \textit{\textbf{Полякова М.В.}}, докт. техн. наук, проф., професор кафедри\\ \> прикладної математики та інформаційних технологій\\ \> НУ "Одеська Політехніка".\\ 
		%			 \> \textit{\textbf{Батюк А.Є.}}, канд.техн.наук, доцент кафедри АСУ\\ \> Інституту комп'ютерних наук та інформаційних технологій\\ \> НУ "Львівська політехніка". \\
		%			 \> \textit{\textbf{Батюк А.Є.}}, канд.техн.наук, доцент кафедри \\ \> Одеського національного університету імені І. І. Мечникова.\\ \> НУ "Львівська політехніка". 
		%		\end{tabbing}
	%		\begin{tabbing}
		%				Відповідальний \= \hspace{3cm}\= \kill
		%				Відповідальний \> \\
		%				редактор: \> \textit{Тарасенко О.В.}, кандидат техн. наук, доцент
		%		\end{tabbing}
	
	
	
	\bigskip
	\medskip
	
	\centering {\normalsize \textit{Рекомендовано до видання Методичною радою ОНУ імені І.І. Мечникова\\
			(протокол № \underline{\hspace{1cm}} від \underline{\hspace{1cm}}.\underline{\hspace{1cm}}.2026 р.) } } \\
	%		за поданням Вченої ради факультету математики, фізики\\ та інформаційних технологій\\
	%		(протокол № \underline{\hspace{1cm}} від \underline{\hspace{1cm}}.\underline{\hspace{1cm}}.2024 р.)}

\bigskip
\medskip

\hspace{1.25cm} \raggedright \textbf{Платонова Є.В., Страхов Є.М.}\\ 
\textbf{XXX} \hspace{0.25cm}\justifying \indent Основи програмування на мові Python. Структурне програмування [Електронний ресурс]: конспект лекцій для здобув. першого (бакалавр) рівня вищ. освіти спец. E7 Математика/ Є.В. Платонова, Є.М. Страхов. --Електронні текстові дані (1 файл: 364 Кб). Одеса: Одес. нац. ун-т ім. І. І. Мечникова, 2026. -- 55 с.

ISBN   XXX-XXX-XXX-XXX-X 

\bigskip
%\parbox{\textwidth}
\justifying 
{\textit {\small Конспект лекцій є введенням у базові принципи програмування мовою Python та систематизацією основних алгоритмічних конструкцій. У матеріалах конспекту розглядаються три базові конструкції управління виконанням програми, надаються приклади реалізації алгоритмів у Python. \\
	Конспект лекцій розроблений для студентів першого курсу спеціальності E7 «Математика» та може бути використаний для самостійного навчання. Матеріал спрямований на формування 
	алгоритмічного мислення та набуття практичних навичок програмування мовою Python.}

\raggedleft	\textbf{УДК 004.4:004.421 }

\small
\begin{tabbing}
	ISBN   XXX-XXX-XXX-XXX-X \= \hspace{4cm} \= ©  \= Платонова Є.В., \= ©  \= Страхов Є.М., 2026 \\
	\>   \hspace{3cm} ©  Одеський національний університет\\  \> \> \> імені І. І. Мечникова, 2026
\end{tabbing}
%	ISBN   XXX-XXX-XXX-XXX-X                                                     ©  Платонова Є. В., 2026
%																										©  Одеський національний університет	імені І. І. Мечникова, 2025
\end{titlepage}

\linespread{1.25}

\KOMAoptions{fontsize=14pt} 


	
		{\large
		\section*{Лабораторна робота 3: Умовні та логічні оператори оператори в Python}
	}
	
	\LabHeader{3}{1}
	
	{\large
		\begin{enumerate}
			\item \textbf{Наступний день.} Дата деякого дня визначається трьома натуральними числами: $d$ (день), $m$ (місяць), $y$ (рік). Визначити дату наступного дня. Розглянути випадки високосного і невисокосного років.
			\item \textbf{Квадрат.} Точки $A$, $B$, $C$, $D$ на площині задаються своїми координатами. Визначити, чи є чотирикутник $ABCD$ квадратом.
		\end{enumerate}
	}
	
	
	\vspace{1cm}
	\LabHeader{3}{2}
	
	{\large
		\begin{enumerate}
			\item \textbf{Упертий Біллі.} Юний Білл Гейтс не любив уроки граматики в школі і якось вирішив втекти з уроку через вікно. Кватирка має форму прямокутника розміром $m\times n$ сантиметрів. Визначити, чи зможе голова Білла Гейтса пролізти у кватирку, якщо голова є ідеальною кулею діаметром $d$ сантиметрів.
			\item \textbf{Прямокутник.} Точки $A$ і $C$ на площині задані своїми координатами і є про\-ти\-леж\-ни\-ми вершинами прямокутника, сторони якого паралельні осям координат. Визначити периметр та площу частини прямокутника, розташованої в першій координатній чверті.
		\end{enumerate}
	}
	
	
	
	\LabHeader{3}{3}
	
	{\large
		\begin{enumerate}
			%\item {\bfseries Игра <<Угадай число 1.0>>.} Компьютер <<загадывает>> случайное целое число от 0 до 9. Игроку предоставляется три попытки угадать это число. В ходе <<диалога>> компьютер не даёт подсказок, а только проверяет пра\-виль\-ность догадки. В случае успеха игроку выдаётся сообщение о выигрыше, а после третьей неудачной попытки "--- о проигрыше.
			\item \textbf{Великий комбінатор.} Новому директору фірми <<Роги і копита>> знадобилася нова шафа. Допоможіть Остапу визначити, чи пройде шафа у дверний проріз, якщо шафа має форму прямокутного паралелепіпеда зі сторонами $x$, $y$, $z$, а дверний проріз має прямокутну форму $a\times b$.
			\item \textbf{Перетин відрізків.} На площині задані чотири точки $A$, $B$, $C$, $D$ такі, що кожен із відрізків $AB$ і $CD$ паралельний одній з осей координат. Перевірити, чи перетинаються відрізки $AB$ та $CD$.
		\end{enumerate}
	}
	\vspace{0.5cm}
	
	\LabHeader{3}{4}
	
	{\large
		\begin{enumerate}
			\item \textbf{Хмарочос.} У хмарочосі $N$ квартир і один під'їзд. На кожному поверсі три квартири. Ліфт може зупинятися лише на непарних поверхах. Пасажир входить до кабіни ліфта і набирає номер необхідної йому квартири $M$. Визначити, на який поверх ліфт має доставити пасажира.
			\item \textbf{Об'єднання прямокутників.} Два прямокутники на площині задані таким чином:
			$$
			\begin{array}{c}
				P_1 = \left\{(x, y):\quad a_1\leqslant x\leqslant a_2,\quad b_1\leqslant y\leqslant b_2\right\},\\ 
				P_2 = \left\{(x, y):\quad c_1\leqslant x\leqslant c_2,\quad d_1\leqslant y\leqslant d_2\right\}.
			\end{array}
			$$
			Знайти площу {\itshape об'єднання} цих прямокутників.
		\end{enumerate}
	}
	

	
	\LabHeader{3}{5}
	
	{\large
		\begin{enumerate}
			\item \textbf{Щасливий квиток.} За заданим шестизначним номером квитка визначити, чи він є <<щасливим>>. Квиток вважається <<щасливим>>, якщо сума перших трьох його цифр дорівнює сумі трьох останніх, але сума не дорівнює 13.
			\item \textbf{Паралельність.} Точки $A$, $B$, $C$, $D$ на площині задані своїми координатами. Чи паралельні прямі $AB$ і $CD$?
		\end{enumerate}
	}
	
	
	
	\LabHeader{3}{6}
	
	{\large
		\begin{enumerate}
			\item \textbf{Касовий апарат.} Скласти програму, яка запитує у користувача вартість покупки і суму грошей, внесену покупцем. Якщо здача не потрібна, надрукувати на екрані <<Дякую>>; якщо внесено більше грошей, ніж потрібно, надрукувати <<Візьміть здачу... грн>> (із зазначенням потрібної суми); якщо ж грошей недостатньо, надрукувати <<Для покупки внесіть ще ... грн>> також із зазначенням необхідної суми.
			\item \textbf{Трикутник і кола.} На площині задані координати вершин трикутника $ABC$ і два кола з центрами $\left(a_1, b_1\right)$, $\left(a_2, b_2\right)$ і радіусами $ r_1$, $r_2$ відповідно. Чи лежить цей трикутник повністю в області перетину цих кіл?
		\end{enumerate}
	}
	
	
	
	\LabHeader{3}{7}
	
	{\large
		\begin{enumerate}
			\item \textbf{Турист.} Турист рухався $t_1$ годин зі швидкістю $v_1$, потім $t_2$ годин "--- зі швидкістю $v_2$ і $t_3$ годин "--- зі швидкістю $v_3$. Визначити, за який час він подолав половину маршруту.
			\item \textbf{Відрізок.} Точки $A$ та $B$ на площині задані своїми координатами. Чи правильно, що відрізок $AB$ не паралельний жодній з координатних осей?
		\end{enumerate}
	}
	
	
	
	\LabHeader{3}{8}
	
	{\large
		\begin{enumerate}
			%\item \textbf{Счастливый билет.} По заданному шестизначному номеру билета определить, является ли он <<счастливым>>. Билет считается <<счастливым>>, если сумма первых трёх его цифр равна сумме трёх последних, при этом сумма не должна равняться 13.
			\item \textbf{Аліса та Боб.} Дані дати народження Аліси та Боба. Кожна з дат складається з трьох чисел: день, місяць та рік. Визначити, хто з них старший.
			\item \textbf{Прямокутний трикутник.} Дано три довільні числа $a$, $b$, $c$. Визначити, чи є трикутник зі сторонами $a$, $b$ і $c$ прямокутним. Якщо такого трикутника не існує, вивести повідомлення. Програма має працювати коректно для будь-якого набору із трьох чисел.
		\end{enumerate}
	}
	
	
	
	\LabHeader{3}{9}
	
	{\large
		\begin{enumerate}
			\item \textbf{Чемпіонат.} Відомі результати двох спортивних команд за підсумками чемпіонату: кіль\-кість очок, число забитих м'ячів, число пропущених м'ячів. Ви\-зна\-чи\-ти, яка з команд вище в~турнірній таблиці. При рівності очок спочатку враховується різниця забитих і пропущених м'ячів, а потім кількість забитих м'ячів. Випадок, коли результати команд ідентичні, не розглядати.
			\item \textbf{Трикутник.} Дані координати вершин трикутника на площині. Визначити, чи паралельна одна з його сторін осі абсцис.
		\end{enumerate}
	}
	
	\vspace{1cm}
	
	\LabHeader{3}{10}
	
	{\large
		\begin{enumerate}
			\item \textbf{Попередній день.} Дата деякого дня визначається трьома натуральними числами: $d$ (день), $m$ (місяць), $y$ (рік). Визначити дату попереднього дня. Рік не є високосним.
			\item \textbf{Перетин прямокутників.} Два прямокутники на площині задані таким чином:
			$$
			\begin{array}{c}
				P_1 = \left\{(x, y):\quad a_1\leqslant x\leqslant a_2,\quad b_1\leqslant y\leqslant b_2\right\},\\ 
				P_2 = \left\{(x, y):\quad c_1\leqslant x\leqslant c_2,\quad d_1\leqslant y\leqslant d_2\right\}.
			\end{array}
			$$
			Знайти площу {\itshape перетину} цих прямокутників. Якщо перетин порожній, вивести~0.
			
		\end{enumerate}
	}
	
	\vspace{1cm}
	
	\LabHeader{3}{11}
	
	{\large
		\begin{enumerate}
			%\item \textbf{Телефонная компания.} Услуги телефонной сети оплачиваются по следующему правилу: за разговоры до $A$ минут в~месяц включительно платится фиксированная сумма $B$ грн, а разговоры сверх установленной нормы оплачиваются по тарифу $C$ грн за минуту. Пользователь вводит количество минут разговоров за прошедший месяц. Вычислить плату за пользование телефоном.
			\item \textbf{Дати.} Дана дата з трьох чисел (день, місяць та рік). Вивести <<yes>>, якщо така дата існує (наприклад, 12.02.1999 "--- <<yes>>, 22.13.2001 "--- <<no>>). Розглянути випадки високосного та невисокосного років.
			\item \textbf{Перпендикулярність.} Дві прямі на площині задані своїми рівняннями $a_1 x+b_1 y +c_1=0$ і $a_2 x+b_2 y +c_2=0$. Визначити, чи перпендикулярні вони.
		\end{enumerate}
	}
	
	
	\vspace{1cm}
	\LabHeader{3}{12}
	
	{\large
		\begin{enumerate}
			%\item {\bfseries Игра <<Угадай число 1.0>>.} Компьютер <<загадывает>> случайное целое число от 0 до 9. Игроку предоставляется три попытки угадать это число. В ходе <<диалога>> компьютер не даёт подсказок, а только проверяет пра\-виль\-ность догадки. В случае успеха игроку выдаётся сообщение о выигрыше, а после третьей неудачной попытки "--- о проигрыше.
			\item \textbf{Робот.} Робот може переміщатися в чотирьох напрямках (<<11>> "--- північ, <<12>> "--- захід, <<13>> "--- південь, <<14>> "--- схід) і приймати три цифрові команди: 0 "--- продовжувати рух, 1 "--- поворот наліво, $-1$ "--- поворот направо. Дано число (11, 12, 13 або 14) "--- початковий напрямок робота і ціле число $N$ (0, 1 або $-1$) "--- послана йому команда. Вивести напрямок робота після виконання отриманої команди (тобто північ, захід, південь або схід).
			\item \textbf{Прямокутник.} Точки $A$ і $C$ на площині задані своїми координатами і є про\-ти\-леж\-ни\-ми вершинами прямокутника, сторони якого паралельні осям координат. Визначити, на скільки частин осі координат розбивають прямокутник. Знайти площу кожної частини.
		\end{enumerate}
	}
	
	\vspace{1cm}
	
	\LabHeader{3}{13}
	
	{\large
		\begin{enumerate}
			\item \textbf{Земельна ділянка.} Чи можна на прямокутній ділянці землі розміром $a\times b$ метрів розмістити дві будівлі з прямокутними основами розмірами $p\times q$ і $r\times s$ метрів? Будівлі можна розташовувати лише паралельно сторонам ділянки.
			\item \textbf{Прямокутний трикутник.} Точки $A$, $B$, $C$ на площині задані своїми координатами. Визначити, чи є трикутник $ABC$ прямокутним. Розглянути всі можливі випадки розташування точок (збіг будь-яких двох точок; всі точки лежать на одній прямій). У разі виродженого трикутника вивести відповідне повідомлення.
		\end{enumerate}
	}
	
	\vspace{1cm}
	
	\LabHeader{3}{14}
	
	{\large
		\begin{enumerate}
			\item \textbf{День тижня.} Дано натуральне число $k$ ($1\leqslant k\leqslant 365$). Визначити, яким днем тижня буде $k$-й день року, якщо 1 січня цього року "--- вівторок.
			\item \textbf{Ромб.} Точки $A$, $B$, $C$, $D$ на площині задаються своїми координатами. Визначити, чи є чотирикутник $ABCD$ ромбом.
		\end{enumerate}
	}
	
	\vspace{1cm}
	
	\LabHeader{3}{15}
	
	{\large
		\begin{enumerate}
			\item \textbf{Вік.} Дано натуральне число $n\leqslant 100$ "--- вік людини. Вивести фразу виду: <<Мені 21 рік>>, <<Мені 12 років>>, <<Мені 32 роки>>.
			\item \textbf{Два прямокутники.} Два прямокутники на площині задані таким чином:
			$$
			\begin{array}{c}
				P_1 = \left\{(x, y):\quad a_1\leqslant x\leqslant a_2,\quad b_1\leqslant y\leqslant b_2\right\},\\ 
				P_2 = \left\{(x, y):\quad c_1\leqslant x\leqslant c_2,\quad d_1\leqslant y\leqslant d_2\right\}.
			\end{array}
			$$
			Чи перетинаються вони? (Випадок однієї спільної точки також вважати перетином.)
		\end{enumerate}
	}
	
	\vspace{1cm}
	
	
	% === ЗМІСТ ===
	\tableofcontents
	

	
	% === СПИСОК ЛІТЕРАТУРИ ===
	% Додайте бібліографічні посилання тут, використовуючи, наприклад, пакет biblatex/natbib
	% \bibliographystyle{plainnat}
	% \bibliography{references} 
	
	
\end{document}